\chapter{Standard Library}
\label{ch:stdlib}

The \haira{} standard library provides essential functionality for common programming tasks. All modules must be explicitly imported.

\section{Core Modules}

\begin{table}[h]
\centering
\begin{tabular}{ll}
\toprule
\textbf{Module} & \textbf{Description} \\
\midrule
\code{io} & Input/output operations \\
\code{string} & String manipulation \\
\code{array} & Array operations \\
\code{map} & Map/dictionary operations \\
\code{math} & Mathematical functions \\
\code{conv} & Type conversions \\
\bottomrule
\end{tabular}
\caption{Core modules}
\end{table}

\section{io Module}

Basic input/output operations.

\begin{lstlisting}
import "io"

// Output
io.print("Hello")           // No newline
io.println("Hello")         // With newline
io.printf("Value: %d\n", 42)  // Formatted

// Input
line = io.readln()          // Read line from stdin

// Errors
io.eprintln("Error!")       // Print to stderr
\end{lstlisting}

\subsection{io Functions}

\begin{lstlisting}
fn print(s: string)
fn println(s: string)
fn printf(format: string, args: ...any)
fn readln() -> string
fn eprintln(s: string)
fn eprintf(format: string, args: ...any)
\end{lstlisting}

\section{string Module}

String manipulation functions.

\begin{lstlisting}
import "string"

s = "  Hello, World!  "

// Basic operations
string.len(s)                    // 18
string.is_empty("")              // true

// Trimming
string.trim(s)                   // "Hello, World!"
string.trim_left(s)              // "Hello, World!  "
string.trim_right(s)             // "  Hello, World!"

// Case
string.to_upper(s)               // "  HELLO, WORLD!  "
string.to_lower(s)               // "  hello, world!  "

// Search
string.contains(s, "World")      // true
string.starts_with(s, "  H")     // true
string.ends_with(s, "!  ")       // true
string.index_of(s, "o")          // 4
string.last_index_of(s, "o")     // 8

// Split/Join
string.split("a,b,c", ",")       // ["a", "b", "c"]
string.join(["a", "b"], "-")     // "a-b"

// Substring
string.substring(s, 2, 7)        // "Hello"
string.char_at(s, 2)             // "H"

// Replace
string.replace(s, "World", "Haira")  // "  Hello, Haira!  "
string.replace_all(s, " ", "_")      // "__Hello,_World!__"

// Repeat
string.repeat("ab", 3)           // "ababab"
\end{lstlisting}

\subsection{string Functions}

\begin{lstlisting}
fn len(s: string) -> int
fn is_empty(s: string) -> bool
fn trim(s: string) -> string
fn trim_left(s: string) -> string
fn trim_right(s: string) -> string
fn to_upper(s: string) -> string
fn to_lower(s: string) -> string
fn contains(s: string, sub: string) -> bool
fn starts_with(s: string, prefix: string) -> bool
fn ends_with(s: string, suffix: string) -> bool
fn index_of(s: string, sub: string) -> int
fn last_index_of(s: string, sub: string) -> int
fn split(s: string, sep: string) -> []string
fn join(parts: []string, sep: string) -> string
fn substring(s: string, start: int, end: int) -> string
fn char_at(s: string, index: int) -> string
fn replace(s: string, old: string, new: string) -> string
fn replace_all(s: string, old: string, new: string) -> string
fn repeat(s: string, n: int) -> string
\end{lstlisting}

\section{array Module}

Array operations.

\begin{lstlisting}
import "array"

arr = [1, 2, 3, 4, 5]

// Basic
array.len(arr)                   // 5
array.is_empty([])               // true

// Access
array.first(arr)                 // 1 (or nil if empty)
array.last(arr)                  // 5 (or nil if empty)
array.get(arr, 2)                // 3 (or nil if out of bounds)

// Modification (returns new array)
array.push(arr, 6)               // [1, 2, 3, 4, 5, 6]
array.pop(arr)                   // ([1, 2, 3, 4], 5)
array.insert(arr, 0, 0)          // [0, 1, 2, 3, 4, 5]
array.remove(arr, 2)             // [1, 2, 4, 5]

// Slicing
array.slice(arr, 1, 4)           // [2, 3, 4]
array.take(arr, 3)               // [1, 2, 3]
array.drop(arr, 2)               // [3, 4, 5]

// Search
array.contains(arr, 3)           // true
array.index_of(arr, 3)           // 2
array.find(arr, fn(x) { x > 3 }) // 4

// Transform
array.map(arr, fn(x) { x * 2 })  // [2, 4, 6, 8, 10]
array.filter(arr, fn(x) { x > 2 })  // [3, 4, 5]
array.reduce(arr, 0, fn(acc, x) { acc + x })  // 15

// Sort
array.sort(arr)                  // [1, 2, 3, 4, 5]
array.sort_by(arr, fn(a, b) { b - a })  // [5, 4, 3, 2, 1]

// Other
array.reverse(arr)               // [5, 4, 3, 2, 1]
array.concat(arr, [6, 7])        // [1, 2, 3, 4, 5, 6, 7]
array.flatten([[1, 2], [3, 4]])  // [1, 2, 3, 4]
array.unique([1, 2, 2, 3])       // [1, 2, 3]
\end{lstlisting}

\section{map Module}

Map/dictionary operations.

\begin{lstlisting}
import "map"

m = ["a": 1, "b": 2, "c": 3]

// Basic
map.len(m)                       // 3
map.is_empty(m)                  // false

// Access
map.get(m, "a")                  // 1 (or nil if not found)
map.has(m, "a")                  // true

// Modification (returns new map)
map.set(m, "d", 4)               // ["a": 1, "b": 2, "c": 3, "d": 4]
map.remove(m, "b")               // ["a": 1, "c": 3]

// Iteration helpers
map.keys(m)                      // ["a", "b", "c"]
map.values(m)                    // [1, 2, 3]
map.entries(m)                   // [("a", 1), ("b", 2), ("c", 3)]

// Transform
map.map_values(m, fn(v) { v * 2 })  // ["a": 2, "b": 4, "c": 6]
map.filter(m, fn(k, v) { v > 1 })   // ["b": 2, "c": 3]

// Merge
map.merge(m, ["d": 4])           // ["a": 1, "b": 2, "c": 3, "d": 4]
\end{lstlisting}

\section{math Module}

Mathematical functions.

\begin{lstlisting}
import "math"

// Constants
math.PI                          // 3.14159265358979
math.E                           // 2.71828182845905

// Basic
math.abs(-5)                     // 5
math.min(3, 7)                   // 3
math.max(3, 7)                   // 7
math.clamp(5, 0, 10)             // 5
math.clamp(-5, 0, 10)            // 0

// Rounding
math.floor(3.7)                  // 3.0
math.ceil(3.2)                   // 4.0
math.round(3.5)                  // 4.0
math.trunc(3.7)                  // 3.0

// Power/Root
math.pow(2.0, 10.0)              // 1024.0
math.sqrt(16.0)                  // 4.0
math.cbrt(27.0)                  // 3.0

// Exponential/Logarithm
math.exp(1.0)                    // 2.718...
math.log(math.E)                 // 1.0
math.log10(100.0)                // 2.0
math.log2(8.0)                   // 3.0

// Trigonometry
math.sin(0.0)                    // 0.0
math.cos(0.0)                    // 1.0
math.tan(0.0)                    // 0.0
math.asin(0.0)                   // 0.0
math.acos(1.0)                   // 0.0
math.atan(0.0)                   // 0.0
math.atan2(1.0, 1.0)             // 0.785... (PI/4)

// Random
math.random()                    // Random float 0.0..1.0
math.random_int(1, 100)          // Random int 1..100
\end{lstlisting}

\section{conv Module}

Type conversion utilities.

\begin{lstlisting}
import "conv"

// To string
conv.int_to_string(42)           // "42"
conv.float_to_string(3.14)       // "3.14"
conv.bool_to_string(true)        // "true"

// From string
conv.string_to_int("42")         // (42, nil)
conv.string_to_int("abc")        // (0, Error)
conv.string_to_float("3.14")     // (3.14, nil)
conv.string_to_bool("true")      // (true, nil)

// Number conversions
conv.int_to_float(42)            // 42.0
conv.float_to_int(3.7)           // 3 (truncates)

// Base conversions
conv.int_to_hex(255)             // "ff"
conv.int_to_binary(10)           // "1010"
conv.int_to_octal(8)             // "10"
conv.hex_to_int("ff")            // (255, nil)
\end{lstlisting}

\section{Extended Modules}

\subsection{fs Module}

File system operations.

\begin{lstlisting}
import "fs"

// Read/Write
content, err = fs.read_file("file.txt")
err = fs.write_file("file.txt", "content")
err = fs.append_file("file.txt", "more")

// File operations
exists = fs.exists("file.txt")
err = fs.remove("file.txt")
err = fs.rename("old.txt", "new.txt")
err = fs.copy("src.txt", "dst.txt")

// Directory operations
err = fs.mkdir("dir")
err = fs.mkdir_all("path/to/dir")
entries, err = fs.read_dir("dir")
err = fs.remove_all("dir")

// File info
info, err = fs.stat("file.txt")
info.size                        // Size in bytes
info.is_dir                      // Is directory?
info.modified                    // Modification time
\end{lstlisting}

\subsection{os Module}

Operating system interface.

\begin{lstlisting}
import "os"

// Environment
os.getenv("HOME")                // Get env variable
os.setenv("KEY", "value")        // Set env variable
os.environ()                     // All env variables

// Process
os.args                          // Command line arguments
os.exit(0)                       // Exit with code
os.getpid()                      // Process ID

// Execution
output, err = os.exec("ls", ["-la"])
\end{lstlisting}

\subsection{time Module}

Time and duration.

\begin{lstlisting}
import "time"

// Current time
now = time.now()
now.year                         // 2024
now.month                        // 1-12
now.day                          // 1-31
now.hour                         // 0-23
now.minute                       // 0-59
now.second                       // 0-59

// Formatting
time.format(now, "2006-01-02")   // "2024-03-15"

// Parsing
t, err = time.parse("2024-03-15", "2006-01-02")

// Duration
time.sleep(1000)                 // Sleep 1 second (ms)
duration = time.since(start)     // Time since start

// Timers
timer = time.after(5000)         // Channel that fires after 5s
ticker = time.tick(1000)         // Channel that fires every 1s
\end{lstlisting}

\subsection{json Module}

JSON encoding/decoding.

\begin{lstlisting}
import "json"

struct User {
    name: string
    age: int
}

// Encoding
user = User{name: "Alice", age: 30}
data, err = json.encode(user)    // '{"name":"Alice","age":30}'

// Decoding
user, err = json.decode<User>(data)

// Pretty printing
data, err = json.encode_pretty(user)

// Raw JSON
value, err = json.parse('{"key": "value"}')
value["key"]                     // "value"
\end{lstlisting}

\subsection{http Module}

HTTP client and server.

\begin{lstlisting}
import "http"

// GET request
response, err = http.get("https://api.example.com/users")
response.status                  // 200
response.body                    // Response body

// POST request
response, err = http.post(
    "https://api.example.com/users",
    ["Content-Type": "application/json"],
    '{"name": "Alice"}'
)

// Server
server = http.Server{port: 8080}

server.handle("/", fn(req: http.Request) -> http.Response {
    return http.Response{
        status: 200,
        body: "Hello, World!"
    }
})

server.listen()
\end{lstlisting}

\subsection{regex Module}

Regular expressions.

\begin{lstlisting}
import "regex"

// Compile pattern
pattern, err = regex.compile(r"\d+")

// Match
regex.is_match(pattern, "abc123")  // true

// Find
match = regex.find(pattern, "abc123def")  // "123"
matches = regex.find_all(pattern, "a1b2c3")  // ["1", "2", "3"]

// Replace
regex.replace(pattern, "a1b2", "X")     // "aXb2"
regex.replace_all(pattern, "a1b2", "X") // "aXbX"

// Capture groups
pattern, _ = regex.compile(r"(\w+)@(\w+)")
captures = regex.captures(pattern, "alice@example")
// captures[0] = "alice@example"
// captures[1] = "alice"
// captures[2] = "example"
\end{lstlisting}

\subsection{crypto Module}

Cryptographic functions.

\begin{lstlisting}
import "crypto"

// Hashing
crypto.md5("hello")              // Hash as hex string
crypto.sha256("hello")
crypto.sha512("hello")

// HMAC
crypto.hmac_sha256("message", "key")

// Random
crypto.random_bytes(32)          // 32 random bytes

// Encoding
crypto.base64_encode(bytes)
crypto.base64_decode(string)
crypto.hex_encode(bytes)
crypto.hex_decode(string)
\end{lstlisting}

\subsection{net Module}

Low-level networking.

\begin{lstlisting}
import "net"

// TCP client
conn, err = net.dial("tcp", "example.com:80")
conn.write("GET / HTTP/1.1\r\n\r\n")
data, err = conn.read(1024)
conn.close()

// TCP server
listener, err = net.listen("tcp", ":8080")
while true {
    conn, err = listener.accept()
    spawn handle_connection(conn)
}
\end{lstlisting}

\section{Testing Module}

\begin{lstlisting}
import "testing"

fn test_addition(t: testing.T) {
    result = add(2, 3)
    t.assert_eq(result, 5)
}

fn test_division(t: testing.T) {
    result, err = divide(10, 2)
    t.assert_nil(err)
    t.assert_eq(result, 5)

    _, err = divide(10, 0)
    t.assert_not_nil(err)
}
\end{lstlisting}

Run tests with:

\begin{lstlisting}[language=bash,numbers=none]
$ haira test
\end{lstlisting}
